% Gemeinsames A4-Querformat mit A5-Text und karierter Notizhälfte.
\usepackage{fontspec}
\setmainfont{Latin Modern Roman}
\usepackage[ngerman]{babel}
\usepackage{amsmath,amssymb,amsthm,mathtools}
\usepackage{booktabs,tabularx,array}
\usepackage{tikz}
\usetikzlibrary{arrows.meta}
\usepackage[a4paper,landscape,top=18mm,bottom=16mm,left=12mm,right=160.5mm]{geometry}
\usepackage{enumitem}
\makeatletter
\renewcommand\normalsize{\@setfontsize\normalsize{10.5}{13}}
\makeatother
\normalsize
\setlist{nosep,leftmargin=1.5em,topsep=.35em}
\pagestyle{empty}
\usepackage{microtype}
\usepackage{needspace}
\usepackage{etoolbox}
\usepackage[hidelinks]{hyperref}
\definecolor{pascalblue}{RGB}{22,89,112}
\definecolor{pascalpale}{RGB}{231,241,244}
\newcommand{\N}{\mathbb N}
\newcommand{\GuideRef}[2]{\href{../main-results/_B10-main-results.pdf\#formula#1}{#2}}
\newcommand{\NewtonLink}{\href{../newton/_B21-newton-reading.pdf\#newton.reading}{Lesefassung zur Newtondarstellung und zu endlichen Differenzen}}
\newtheorem*{satz}{Satz}
\newtheorem*{definition}{Definition}
\theoremstyle{remark}
\newtheorem*{beispiel}{Beispiel}
\newcommand{\Beweisidee}{\par\smallskip\textbf{Beweisidee.}\ }
\newcommand{\PascalDreieck}[1]{%
\begin{tikzpicture}[x=0.62cm,y=0.52cm,
 entry/.style={minimum width=6mm,minimum height=5mm,inner sep=1pt,font=\small},
 binhighlight/.style={fill=pascalpale,text=pascalblue,rounded corners=2pt,font=\small\bfseries}]
 \foreach \n/\row in {0/{1},1/{1,1},2/{1,2,1},3/{1,3,3,1},4/{1,4,6,4,1},5/{1,5,10,10,5,1},6/{1,6,15,20,15,6,1}} {
  \node[anchor=east,font=\small,text=black!65] at (-7,-\n) {$n=\n$};
  \foreach \v [count=\k from 0] in \row {
   \ifnum#1=1
    \ifnum\n=4
     \ifnum\k=2 \node[entry,binhighlight] at ({2*\k-\n},-\n) {\v};
     \else \node[entry] at ({2*\k-\n},-\n) {\v};\fi
    \else \node[entry] at ({2*\k-\n},-\n) {\v};\fi
   \else
    \pgfmathtruncatemacro{\selected}{(\k==2 && \n>=2 && \n<=5) || (\n==6 && \k==3)}
    \ifnum\selected=1 \node[entry,binhighlight] at ({2*\k-\n},-\n) {\v};
    \else \node[entry] at ({2*\k-\n},-\n) {\v};\fi
   \fi
  }
 }
 \ifnum#1=1
  \draw[-{Stealth[length=1.5mm]},pascalblue,shorten >=3mm,shorten <=3mm] (-1,-3) -- (0,-4);
  \draw[-{Stealth[length=1.5mm]},pascalblue,shorten >=3mm,shorten <=3mm] (1,-3) -- (0,-4);
 \fi
\end{tikzpicture}}
\setlength{\parindent}{0pt}
\setlength{\parskip}{.42em plus .08em minus .08em}
\setlength{\emergencystretch}{1.5em}
\setcounter{tocdepth}{1}
\newlength{\BinWorkHeight}
\newcommand{\ExerciseSpace}[1]{%
 \par\smallskip\begingroup
 \setlength{\BinWorkHeight}{\dimexpr\pagegoal-\pagetotal-2\baselineskip\relax}%
 \ifdim\BinWorkHeight>#1mm\setlength{\BinWorkHeight}{#1mm}\fi
 \ifdim\BinWorkHeight>12mm
  \vfill{\small\sffamily Deine Rechnung oder Skizze}\par\nobreak
  \begin{tikzpicture}
   \draw[step=5mm,black!22,line width=.2pt] (0,0) grid (120mm,\BinWorkHeight);
  \end{tikzpicture}
 \fi\endgroup}
\newcommand{\Exercise}[2]{%
 \clearpage
 \ifnum#1=1
  \section{Übungsblätter}\label{bin:exercises}
  Arbeite zuerst ohne die Lösungen ab Seite~\pageref{bin:solutions}.
  Nutze das Karofeld rechts und den freien Raum unter der Aufgabe.
 \fi
 \subsection*{Aufgabe #1: #2}\phantomsection\label{bin:exercise-#1}}
\newcommand{\ExtraWork}[2]{%
 \clearpage\subsection*{Arbeitsraum zu Aufgabe #1}
 \textbf{Hinweis, falls du ihn brauchst.} #2\par
 \ExerciseSpace{130}}
\AddToHook{shipout/background}{%
 \begin{tikzpicture}[remember picture,overlay]
  \begin{scope}[shift={(current page.north west)},x=1mm,y=-1mm]
   \ifnum\value{page}>1
    \node[anchor=west,font=\sffamily\fontsize{7.5}{9}\selectfont,text=black!55]
       at (12,10) {BAND 10\quad BINOMIALKOEFFIZIENTEN\quad \BinBookletLabel};
    \draw[black!25,line width=.25pt] (12,13)--(136.5,13);
   \fi
   \node[anchor=west,font=\sffamily\fontsize{8.5}{10}\selectfont,text=black!65]
      at (161,13) {NOTIZEN\quad Fragen, Ideen und Rechnungen};
   \begin{scope}[shift={(161,22)}]
    \draw[step=5mm,black!22,line width=.2pt] (0,0) grid (125,170);
   \end{scope}
   \node[font=\sffamily\fontsize{8}{10}\selectfont,text=black!60]
      at (74.25,201) {\BinPageLabel{} \thepage};
   \node[font=\sffamily\fontsize{8}{10}\selectfont,text=black!60]
      at (223.5,201) {Notizen zu \BinPageLabel{} \thepage};
  \end{scope}
 \end{tikzpicture}}
\newcommand{\PascalEltern}{%
 \begin{tikzpicture}[x=1mm,y=1mm,
  parent/.style={draw=pascalblue!50,fill=pascalpale,rounded corners=2pt,inner sep=6pt},
  arr/.style={-{Stealth[length=2mm]},pascalblue,line width=.7pt}]
  \node[parent] (L) at (-24,0) {$\displaystyle\binom n{k-1}$};
  \node[parent] (R) at (24,0) {$\displaystyle\binom nk$};
  \node[parent] (C) at (0,-26) {$\displaystyle\binom{n+1}k$};
  \node[font=\small,anchor=east] at (-43,0) {Zeile $n$};
  \node[font=\small,anchor=east] at (-25,-26) {Zeile $n+1$};
  \draw[arr] (L.south)--(C.north west);
  \draw[arr] (R.south)--(C.north east);
  \node[font=\footnotesize,anchor=south] at (-24,8) {Index $k-1$};
  \node[font=\footnotesize,anchor=south] at (24,8) {Index $k$};
  \node[font=\footnotesize,anchor=north] at (0,-34) {Index $k$ in der neuen Zeile};
 \end{tikzpicture}}
\allowdisplaybreaks[1]
\clubpenalty=10000
\widowpenalty=10000
\BeforeBeginEnvironment{satz}{\Needspace{6\baselineskip}}
\BeforeBeginEnvironment{definition}{\Needspace{6\baselineskip}}
\BeforeBeginEnvironment{beispiel}{\Needspace{5\baselineskip}}
\pretocmd{\section}{\Needspace{9\baselineskip}}{}{}
